\documentstyle[% 


righttag% 


,11pt% 


,amssymb% 


,russian%               remove in Eng. 


,izvru%                 Set izven% in Eng. 


]{amsart} 


 


% 


 


\newcommand{\hateq}{\,\wh=\,} 


\newcommand{\la}{\langle} 


\newcommand{\ra}{\rangle} 


\newcommand{\esssup}{\operatorname{ess\;sup}} 


\newcommand{\essinf}{\operatorname{ess\;inf}} 


\newcommand{\Span}{\operatorname{Span}} 


\newcommand{\tts}[1]{} 


 


\theoremstyle{plain} 


\theorembodyfont{\it} 


\newtheorem{assrt}{\hskip\parindent Утверждение} 


 


%\hoffset=-15mm%                  For Eng. 


 


\setcounter{page}{18} 


\god{2001} 


\nomer{10~(473)} %                        Remove in Eng. 


%\nomer{Vol.\,45, No.\,10}%               Set in Eng. 


%\str{\pageref{begin}--\pageref{end}}%  Set in Eng. 


\udk{517.926} 


 


\author{\it А.В.\,Данеев, В.А.\,Русанов} 


% NB:   ^^ remove in Eng. 


\title{Геометрический подход к решению некоторых\\ обратных 


задач системного анализа} 


\thanks{Работа выполнена при финансовой поддержке Российского фонда 


фундаментальных исследований (гранты \No\,99-01-01279, 


\No\,00-01-00922).} 


 


\date{} 


% \pagestyle{myheadings}%         Set in Eng. 


\pagestyle{plain} %              Remove in Eng. 


\begin{document} 


% \thispagestyle{plain}%          Set in Eng. 


\maketitle 


\label{begin} 


 


 


{\bf Введение.} Аксиоматическое основание теории идентификации сложных 


динамических систем, развитое в \cite{1}, приводит к понятиям 


$I$-процесса и идентификационного пространства --- конструкциям


предложенного в \cite{1} аксиоматического построения. Структура 


идентификационного пространства делает возможным строгое (абстрактное) 


определение свойства идентифицируемости математической модели 


динамической системы, не связывая это понятие с какой-либо процедурой 


или алгоритмом процесса идентификации, а также полное описание 


семейства всех максимальных подмножеств пространства динамических 


систем, идентифицируемых посредством анализа откликов этих систем на 


элементы a~posteriori приобретенного множества входных воздействий. 


 


Конструктивность созданного на этом пути набора общих понятий и методов 


апробирована в \cite{2}--\cite{9} на примерах постановок классических 


задач идентификации и реализации \cite{10} линейной конечномерной 


нестационарной динамической системы управления в классе обыкновенных 


дифференциальных уравнений. Что же касается результатов качественной 


теории идентификации, относящихся к бесконечномерному случаю, то они не 


затрагивают всех вопросов, которые нашли свое законченное решение в 


конечномерном варианте (хотя многие из них решаются уже на этой степени 


общности \cite{11}). 


 


Цель предлагаемой статьи заключается в рассмотрении на строгой 


теоретико-системной основе методологии по идентификации моделей 


линейных динамических объектов с уравнениями состояния в классе 


обыкновенных дифференциальных уравнений в бесконечномерном банаховом 


пространстве. Излагаемые ниже элементы общей теории $I$-процессов в 


общем банаховом пространстве (не обязательно сепарабельном и равномерно 


выпуклом (\cite{12}, с.\,182)) не претендуют на полноту и 


законченность, а только уточняют главные принципы идентификации. 


\medskip 


 


{\bf 1. Вспомогательные построения.} В этом разделе приводятся основные 


определения и обозначения, а также устанавливаются вспомогательные 


утверждения. 


 


Пусть $X$ и $Y$ --- вещественные (не нулевые) линейные пространства, 


элементы которых определяют соответственно {\it вход} и {\it выход} 


математических моделей динамических систем из a~priori заданного 


подсемейства $L(X,Y)\subset Y^X$; здесь $L(X,Y)$ --- семейство всех 


линейных операторов из $X$ в $Y$, $Y^X$ --- семейство всех функций, 


отображающих множество $X$ в $Y$. Ясно, что $L(X,Y)$ и $Y^X$ сами 


являются векторными пространствами относительно обычных операций 


сложения функций и умножения их на действительные числа (при этом 


играет роль лишь наличие структуры векторного пространства в $Y$, а не 


в $X$). 


 


Следуя \cite{1}, \cite{2}, будем говорить, что идентификационное 


пространство --- это упорядоченная пара $(L(X,Y),\frak J)$, состоящая 


из линейного множества $L(X,Y)$ и некоторого семейства $\frak J$ 


подмножеств (бинарных отношений) множества $L(X,Y)\times L(X,Y)$, 


удовлетворяющего (в терминах алгебры отношений) условиям: 


 


--- любой элемент из $\frak J$ содержит диагональ $\Delta_L\subset 


L(X,Y)\times L(X,Y)$; 


 


--- если $V\in\frak J$, то $V^{-1}=V$; 


 


--- если $V\in\frak J$, то $V\circ V\subset V$; 


 


--- если $\frak R\subset\frak J$, то $\cap\frak R\in\frak J$. 


 


Множество $L(X,Y)$ называется пространством {\it линейных 


математических моделей} динамических систем; подмножества из 


$L(X,Y)\times L(X,Y)$, принадлежащие $\frak J$, называются {\it 


идентификационными признаками} в $L(X,Y)$; семейство $\frak J$ 


идентификационных признаков пространства $L(X,Y)$ называется {\it 


идентификацией} на $L(X,Y)$ \cite{1}. 


 


Будем говорить, что пространство $(L(X,Y),\frak J)$ обладает свойством 


{\it идентифицируемости} ({\it идентифицируемо}), если 


$\Delta_L\in\frak J$ \cite{1}. 


 


Функцию $\xi\to\pi_x(\xi):L(X,Y)\to Y$, задаваемую равенством 


$\pi_x(\xi)\hateq\xi(x)$, $x\in X$, назовем {\it вычислением в точке} 


$x$. Используя эту функцию, определим отображение $\pi^*_x:L(X,Y)\times 


L(X,Y)\to Y\times Y$, $\pi^*_x(\sigma,\nu)\hateq 


(\pi_x(\sigma),\pi_x(\nu))$. Обозначим через $\frak J_\Omega$ 


идентификацию пространства $L(X,Y)$ вида $\{V\subset L(X,Y)\times 


L(X,Y):\exists W\subset \Omega$, $V=\bigcap\limits_{x\in 


W}\pi^{*-1}_x[\Delta_Y]\}$, $\Omega\subset X$. В этом положении 


$\Omega$ назовем {\it идентификационным базисом} ($I$-{\it базисом}) 


пространства $(L(X,Y),\frak J_\Omega)$, и если пространство 


$(L(X,Y),\frak J_\Omega)$ идентифицируемо, то будем говорить, что 


$I$-базис $\Omega$ {\it полный}, при этом минимальное кардинальное 


число полного $I$-базиса для $L(X,Y)$ назовем {\it идентификационной 


размерностью} $L(X,Y)$ \cite{1}. 


 


Пусть $\cal E_{\frak J_\Omega}$ --- равномерная топология, порожденная 


на $L(X,Y)$ равномерностью с базой $\frak J_\Omega$. 


 


\begin{lem} 


В топологическом пространстве $(L(X,Y),\cal E_{\frak J_\Omega})$ операция 


сложения непрерывна. 


\end{lem} 


 


\begin{note} 


$(L(X,Y),\cal E_{\frak J_\Omega})$ не является топологическим векторным 


пространством, т.\,к. операция умножения на числа не удовлетворяет 


$\cal E_{\frak J_\Omega}$-непрерывности. 


\end{note} 


 


Следующая теорема устанавливает тождественность задач А и Б из \cite{1} 


в классе $L(X,Y)$; в определении псевдо-идентифицируемости следуем 


\cite{1}. 


 


\begin{thm} 


Идентифицируемость пространства $(L(X,Y),\frak J_\Omega)$ равносильна 


псевдо-иден\-ти\-фицируемости некоторой его точки. 


\end{thm} 


 


Пусть $X$ и $Y$ --- топологические векторные пространства (ТВП), 


удовлетворяющие аксиоме отделимости $T_1$, и $\cal L(X,Y)$ \label{space}


--- пространство всех линейных непрерывных операторов из $X$ в $Y$. Если 


$X$ --- локально выпуклое пространство, что и предполагаем везде далее, 


то для него корректно рассмотрение фундаментального подмножества 


(\cite{13}, с.\,115). 


 


\begin{thm} 


$(\cal L(X,Y),\frak J_\Omega)$ идентифицируемо в том и только том 


случае, если $I$-базис $\Omega$ --- множество, фундаментальное в $X$. 


\end{thm} 


 


В завершение этого раздела опишем в терминах топологии пространства $X$ 


положение, когда полный $I$-базис является конечным. 


 


\begin{thm} 


Идентификационная размерность $\cal L(X,Y)$ конечна тогда и только тогда, 


когда топология в $X$ локально компактна. 


\end{thm} 


 


\begin{cor} 


Если идентификационная размерность семейства $\cal L(X,Y)$ конечна, то 


$X$ --- множество второй категории в себе. 


\end{cor} 


 


{\bf 2. Конструкции идентификационного пространства.} 


В этом разделе будут определены основные положения, связанные с 


построением идентификационных пространств по технологии {\it 


вход--выход} \cite{1} на семействах математических моделей динамических 


систем с уравнениями состояния в классе линейных нестационарных 


обыкновенных дифференциальных уравнений в общем банаховом пространстве. 


 


Везде далее $(X,\|\cdot\|_X)$ и $(Y,\|\cdot\|_Y)$ --- вещественные 


банаховы пространства, $\cal L(Y,X)$ --- банахово пространство (с 


операторной нормой) всех линейных непрерывных операторов, действующих 


из $Y$ в $X$ (аналогично $\cal L(X,X)$), $T\hateq[t_0,t_1]$ --- отрезок 


числовой прямой $R$ с мерой Лебега $\mu$ и $\nu$ --- полная 


положительная мера, абсолютно непрерывная относительно $\mu$ и 


определенная на $\sigma$-алгебре $\frak S_\nu$\; $\nu$-измеримых 


(лебеговски пополненных) подмножеств из $T$. 


 


Обозначим через $L_p(T,\nu,X)$ ($L_p(T,\nu,Y)$), $p\in[1,\infty)$, банахово 


пространство классов эквивалентности всех интегрируемых (по Бохнеру 


\cite{12}) отображений $\varphi:T\to X$ (соответственно $\varphi:T\to 


Y$) с нормой $\|\varphi\|^X_{\nu,p}\hateq\big(\int\limits_T 


\|\varphi(t)\|^p_X \nu(dt)\big)^{1/p}$ (аналогично 


$\|\varphi\|^Y_{\nu,p}$). Через $L_\infty(T,\nu,X)$ 


($L_\infty(T,\nu,Y)$) будем обозначать банахово пространство 


(эквивалентных классов) всех сильно $\nu$-измеримых (\cite{12}, с.\,187) и 


$\nu$-существенно ограниченных функций из $T$ в $X$ ($Y$) с нормой 


$\|\psi\|^X_{\nu,\infty}\hateq \esssup\|\psi(t)\|_X$ 


($\|\psi\|^Y_{\nu,\infty}$). Как обычно, $AC(T,X)$ --- линейное 


многообразие всех абсолютно непрерывных на $T$ функций со значениями в 


$X$. 


 


Выделим класс динамических систем, движение которых в пространстве 


состояний $X$ описывается дифференциальным уравнением вида 


\begin{equation} 


\Dot x(t)=A(t)x(t)+B(t)u(t),\quad t\in T, 


\end{equation} 


где $x(\cdot)\in AC(T,X)$, $A(\cdot)\in L_p(T,\mu,\cal L(X,X))$, 


$B(\cdot)\in L_p(T,\mu,\cal L(Y,X))$ ($p\in[1,\infty]$), $Y$ --- 


банахово пространство управляющих воздействий $u(\cdot)\in 


L_{p'}(T,\mu,Y)$, $p'$ --- число, сопряженное с $p$. Для определенности 


условимся рассматривать класс решений уравнения (1) типа Каратеодори. 


 


Обозначим в целях удобства через $H_{p'}$ и $H_\infty$ векторные 


пространства $L_{p'}(T,\mu,X)\times L_{p'}(T,\mu,Y)$ и 


$L_\infty(T,\mu,X)\times L_\infty(T,\mu,Y)$ с нормами 


$\|(\omega_1,\omega_2)\|_{H_{p'}}\hateq 


[(\|\omega_1\|^X_{\mu,p'})^{p'}+ 


(\|\omega_2\|^Y_{\mu,p'})^{p'}]^{1/p'}$, $\omega_1\in L_{p'}(T,\mu,X)$, 


$\omega_2\in L_{p'}(T,\mu,Y)$ ($p'\in[1,\infty)$) и 


$\|(\omega_1,\omega_2)\|_{H_\infty}\hateq 


\max(\|\omega_1\|^X_{\mu,\infty}, \|\omega_2\|^Y_{\mu,\infty})$, 


$\omega_1\in L_\infty(T,\mu,X)$, $\omega_2\in L_\infty(T,\mu,Y)$, 


соответственно, которые являются банаховыми; если в некотором 


утверждении о пространствах $H_{p'}$ и $H_\infty$ их индексы явно не 


указаны, то это утверждение будем относить к обоим случаям. Как на


с.\,\pageref{space}, $\cal L(H,X)$ --- пространство всех линейных непрерывных 


операторов из $H$ в $X$. 


 


Оператор $\xi:H\to X$, действующий в соответствии с 


\begin{equation} 


\xi(\omega_1,\omega_2)\hateq\int_T(A(t)\omega_1(t)+B(t) 


\omega_2(t))\mu(dt), 


\end{equation} 


является линейным и непрерывным, а значит, $\xi\in \cal L(H,X)$. 


 


\begin{defn} 


Для динамического объекта (1) пару $(A,B)\in L_p(T,\mu,\cal 


L(X,X))\times\linebreak L_p(T,\mu,\cal L(Y,X))$,


где $A(\cdot)$ и $B(\cdot)$ --- 


отображения из (1), будем называть $(A,B)$-моделью, при этом оператор 


$\xi$, действующий в соответствии с (2), назовем $\xi$-моделью, а 


$(A,B)$- и $\xi$-модели, связанные соотношением (2), --- 


$I$-эквивалентными. 


\end{defn} 


 


Следуя \cite{1}, \cite{2}, будем говорить, что в $I$-процессе множество 


$H$ образует {\it пространство входных сигналов}, $X$ --- {\it 


пространство выходных сигналов}, $\cal L(H,X)$ --- {\it пространство 


математических моделей} динамических систем. Поясним необходимость 


введенных понятий. 


 


Классический подход \cite{10} к проблеме идентификации объекта (1) 


обычно имел дело непосредственно с $(A,B)$-моделью. При таком подходе 


построение $I$-процесса по технологии {\it вход--выход}, предложенной в 


\cite{1}, не представляется возможным. С другой стороны, концепция 


$\xi$-модели такую технологию допускает. Далее остаются только детали: 


оператор $\Gamma: L_p(T,\mu,\cal L(X,X))\times L_p(T,\mu,\cal L(Y,X)) 


\to\cal L(H,X)$, действующий в соответствии с (2), линейный и 


взаимно однозначный, следовательно, в соответствии с \cite{1} {\it 


идентифицируемость} переносится с $\cal L(H,X)$ в $L_p(T,\mu,\cal 


L(X,X))\times L_p(T,\mu,\cal L(Y,X))$ при действии $\Gamma$. 


 


\begin{defn} 


Пару $\chi=(x,u)\in AC(T,X)\times L_{p'}(T,\mu,Y)$, $(x,u)\ne0$, 


удовлетворяющую $\mu$-почти всюду соотношению 


(1) для некоторой пары $(A,B)$ операторов $A(\cdot)\in L_p(T,\mu,\cal 


L(X,X))$ и $B(\cdot)\in L_p(T,\mu,\cal L(Y,X))$, назовем порождающим 


(характеристическим) элементом идентификационного процесса 


($I$-процесса) $(A,B)$-модели ($\xi$-модели) динамического объекта (1). 


\end{defn} 


 


Проблема построения $I$-процесса по технологии {\it вход--выход} по 


существу ставится \cite{1} следующим образом: используя отображение 


$t\to\chi(t):T\to X\times Y$, построить $I$-базис $\Omega\subset H$ 


такой, что для любого подмножества $W\subset \Omega$ можно сформировать 


идентификатор $(H,X,\cal L(H,X),\ \bigcap\limits_{\omega\in 


W}\pi^{-1}_\omega\circ\pi_\omega,\ \bigcap\limits_{\omega\in 


W}\pi^{-1}_\omega(\xi(\omega)))$ как элемент $I$-процесса. 


 


\begin{defn} 


Множество $\Lambda^{\#}$ всех вещественных на $T$ функций таких, что 


$\lambda\in\Lambda^{\#}\Rightarrow\lambda\chi\in H$, где $\chi$ 


--- характеристический элемент $I$-процесса динамического объекта (1), 


назовем семейством сигнальных функций элемента $\chi$. (Условимся 


использовать обозначения $\Lambda^{\#}_{p'}$ и $\Lambda^{\#}_\infty$, 


когда того требует уточнение характера пространства входных сигналов 


$H$ ($H_{p'}$ или $H_\infty$).) 


\end{defn} 


 


Как увидим ниже, семейство $\Lambda^{\#}$ допускает явную реализацию в 


виде классов эквивалентностей функций с отождествлением функций из 


$\Lambda^{\#}$, совпадающих почти всюду. 


 


Покажем, что $I$-базисом является подмножество из $H$ вида $\{\omega\in 


H:\omega=\lambda\chi,\ \lambda\in\Lambda\}$, где $\Lambda\subset 


\Lambda^{\#}$. Для этого необходимо показать, что любой 


идентификационный признак $\pi^{-1}_\omega\circ\pi_\omega$ 


($\omega=\lambda\chi$, $\lambda\in\Lambda$) совместно с $\chi$ 


однозначно определяют идентификатор $(H,X,\cal L(H,X)$, 


$\pi^{-1}_\omega\circ\pi_\omega$, $\pi^{-1}_\omega(\xi(\omega)))$ как 


элемент $I$-процесса. При этом определение компонент идентификатора 


связано с тем, что по $\lambda$ и $\chi$ можно вычислить значение 


оператора $\xi$ в точке $\omega=\lambda\chi\in H$. Именно, 


\begin{equation} 


\xi(\omega)=\int_T \lambda(t)\Dot x(t)\mu(dt). 


\end{equation} 


 


Таким образом, характеристический элемент $\chi$ совместно с любым 


подмножеством $\Lambda\subset\Lambda^{\#}$ порождают идентификационный 


базис $\Omega\hateq\{\omega: \omega=\lambda\chi,\ \lambda\in\Lambda\}$, 


а он в свою очередь --- идентификационное пространство $(\cal 


L(H,X),\frak J_{\Omega})$, 


$$ 


\frak J_\Omega\hateq\Big\{V\subset\cal L(H,X)\times \cal L(H,X): 


\exists W\subset\Omega,\ V=\bigcap_{\omega\in W} 


\pi^{*-1}_\omega[\Delta_x]\Big\}. 


$$ 


 


Известно \cite{2}, \cite{5}, что даже в случае конечномерного объекта 


(1) не существует порождающего элемента (соответственно $I$-базиса 


$\Omega$) идентификационного процесса $\xi$-модели такого, что 


пространство $(\cal L(H,X),\frak J_\Omega)$ идентифицируемо. Таким 


образом, центральным местом в общей теории идентификации систем вида 


(1) становится проблема исследования тополого-алгебраической структуры 


максимальных (в упорядочении относительно теоретико-множественного 


включения) идентифицируемых подмножеств из $(\cal L(H,X),\frak 


J_\Omega)$. 


\medskip 


 


{\bf 3. Геометрические свойства семейства сигнальных функций.} В силу 


(3) семейство сигнальных функций $\Lambda^{\#}$ аналогично пространству 


основных (пробных) функций в теории распределений. 


 


\begin{thm} 


{\em a)} $\Lambda^{\#}_{p'}=L_{p'}(T,\nu,R)$, $p'\in[1,\infty);$ {\em б)} 


$\Lambda^{\#}_\infty=L_\infty(T,\nu,R)$, где $\nu$ --- мера, определяемая 


соотношением 


$$ 


\nu(S)\hateq\int_S(\|x(t)\|^{p'}_X+\|u(t)\|^{p'}_Y)\mu(dt),\quad 


S\in\frak S_\nu. 


$$ 


\end{thm} 


 


Это значит, что при заданном $\chi$ семейство $\Lambda^{\#}$ полностью 


характеризуется конструкцией обычного лебегова пространства $L_k$, при 


этом индекс $k=p'$ из формулировки теоремы 4 обладает определенной 


жесткостью, а именно, его нельзя понизить с сохранением включения 


$\lambda\chi\in H$ и нельзя повысить, не теряя свойства максимальности 


$\Lambda^{\#}$. 


 


\begin{note} 


В случае б) индекс $p'$ в соотношении для меры $\nu$, с одной стороны, по 


смыслу не может являться числом, сопряженным для $p$ из (1), а с другой, 


этот индекс не влияет на $\sigma$-алгебру $\frak S_\nu$, поэтому 


выбирается произвольным из $[1,\infty)$. 


\end{note} 


 


{\bf Доказательство.} Так как $\Dot x(\cdot)$ и $u(\cdot)$ сильно 


измеримы, то они $\mu$-почти сепарабельнозначны (\cite{12}, с.\,187). 


Поэтому, не нарушая общности, пространства $X$ и $Y$ можно считать 


сепарабельными. 


 


a) Установим $\mu$-измеримость функций $t\to|\lambda(t)|\,\|x(t)\|_X$ и 


$t\to|\lambda(t)|\,\|u(t)\|_Y$, т.\,к. в этом случае 


доказательство теоремы сводится (детали опускаем, отсылая за ними к 


п.\,3) теоремы 1.4.20 и п.\,1) теорем 1.4.27, 1.4.38 из \cite{14}) к 


установлению следующей цепочки эквивалентностей: 


\begin{multline*} 


\lambda\in\Lambda^{\#}_{p'}\Longleftrightarrow \lambda x\in 


L_{p'}(T,\mu,X)\,\&\, \lambda u\in L_{p'}(T,\mu,Y) 


\Longleftrightarrow \\ 


% 


\Longleftrightarrow|\lambda|\,\|x\|_X,\,|\lambda|\,\|u\|_Y\in 


L_{p'}(T,\mu,R)\Longleftrightarrow |\lambda|^p 


(\|x\|^{p'}_X+\|u\|^{p'}_Y)\in L_1(T,\mu,R)\Longleftrightarrow \\ 


% 


\Longleftrightarrow|\lambda|^{p'}\in L_1(T,\nu,R)\Longleftrightarrow 


\lambda\in L_{p'}(T,\nu,R). 


\end{multline*} 


 


Функция $t\to(\|x(t)\|_X+\|u(t)\|_Y)$\; $\mu$-измерима и, 


следовательно, $\mu$-измеримо множество $S_0=\{t\in 


T:(\|x(t)\|_X+\|u(t)\|_Y)=0\}$. Ясно, что возможны два варианта: 1) 


$\mu(S_0)=0$ и 2) $\mu(S_0)>0$. 


 


1) В силу теоремы 1.4.38 из \cite{14} функция 


$t\to|\lambda(t)|^p\varphi(t)$, где 


$$ 


\varphi(t)\hateq 


\begin{cases} 


(\|x(t)\|^{p'}_X+\|u(t)\|^{p'}_X), & t\in T\setminus S_0;\\ 


\text{const\,}\ne0, & t\in S_0, 


\end{cases} 


$$ 


$\mu$-измерима, откуда следует $\mu$-измеримость $\lambda$ и, 


следовательно, $\mu$-измеримость функций $|\lambda|\,\|x\|_X$ и 


$|\lambda|\|u\|_Y$. 


 


2) Так как $S_0\subset\{t\in T:\|x(t)\|_X=0\}$, то 


$|\lambda|^{p'}\|x\|^{p'}_X(\|x\|^{p'}_X+\|u\|^{p'}_Y)= 


|\lambda|^{p'}\|x\|^{p'}_X\varphi$, где левая часть равенства --- 


$\mu$-измеримая функция (как произведение $\mu$-измеримых функций 


$|\lambda|^{p'}(\|x\|^{p'}_X+\|u\|^{p'}_Y)$ и $\|x\|^{p'}_X$). Умножая 


правую часть равенства на $\mu$-измеримую функцию $1/\varphi$, приходим 


к заключению, что отображение $t\to|\lambda(t)|\,\|x(t)\|_X$\; 


$\mu$-измеримо. Аналогично доказывается $\mu$-измеримость 


$t\to|\lambda(t)|\,\|u(t)\|_Y$. Для доказательства основного результата 


достаточно рассуждений п.\,2), но в 1) вместе с основным утверждением 


была уточнена $\nu$-измеримость $\lambda$ до $\mu$-измеримости. 


 


б) Доказательство следует из цепочки эквивалентностей: 


\begin{multline*} 


\lambda\in\Lambda^{\#}_\infty\Longleftrightarrow \lambda x\in 


L_\infty(T,\mu,X)\,\&\, \lambda u\in L_\infty(T,\mu,Y) 


\Longleftrightarrow \\ 


% 


\Longleftrightarrow |\lambda|\,\|x\|_X,\,|\lambda|\,\|u\|_Y\in 


L^*_1(T,\mu,R)\Longleftrightarrow\\ 


\Longleftrightarrow \forall \psi\ (\psi\in L_1(T,\nu,R)\,\&\, 


\psi\lambda\in L_1(T,\nu,R)) \Longleftrightarrow 


\lambda\in L^*_1(T,\nu,R)\Longleftrightarrow \lambda\in 


L_\infty(T,\nu,R). 


\end{multline*} 


 


\begin{cor} 


a) $\Dot x\in L_p(T,\nu,X)$, $p\in[1,\infty]$; б) $\Lambda^{\#}_{p'} 


\subset L_1(T,\nu_-,R)$, где $\nu_-(S)\hateq\int\limits_S \|\Dot x(t)\|_X 


\mu(dt)$, $S\in\frak S_{\nu_-}$. 


\end{cor} 


 


(Данное следствие весьма важно при построении теории сильных 


$(A,B)$-моделей \cite{3}, \cite{4}, \cite{7}, \cite{9}.) 


 


Пусть в соответствии с теоремой 4 определено семейство $\Lambda^{\#}$. 


В этом случае в качестве семейств сигнальных функций корректны рассмотрения 


подмножеств $\Lambda_k\subset \Lambda^{\#}$ вида $\Lambda_k\hateq 


L_k(T,\nu,R)$, $p'\le k\le\infty$ (здесь и далее мера $\nu$ выбирается 


согласно формулировке теоремы 4). Введем на каждом $\Lambda_k$ 


топологическую структуру, инициированную нормой $\|\cdot\|^R_{\nu,k}$. 


То, что топологические свойства линейных оболочек подмножеств из 


$\Lambda^{\#}$ в некоторых случаях соотносятся с их алгебраической 


размерностью, доказывает 


 


\begin{thm} 


Пусть $\Lambda\subset\Lambda_\infty\subset\Lambda^{\#}$. Тогда 


$\dim\Span\Lambda$ конечна, если и только если $\Span\Lambda$ замкнуто 


в $L_l(T,\nu,R)$, $l\in[1,\infty)$. 


\end{thm} 


 


Необходимость очевидна. Справедливость достаточных условий легко 


установить модификацией теоремы Гротендика (\cite{15}, с.\,133). 


 


Теорема 5 показывает, что иногда по характеру топологии в 


$\Span\Lambda$ можно судить о характере размерности $I$-базиса. 


 


Уточним взаимосвязь между банаховыми пространствами сигнальных функций 


$\Lambda^{\#}$ и выходных сигналов $X$ через действие $\xi$-модели (2) 


в соответствии с соотношением (3). Пусть $\xi_\Lambda:\Lambda^{\#}\to 


X$ --- оператор, задаваемый равенством 


$\xi_\Lambda(\lambda)\hateq\xi(\lambda\chi)$, которое корректно в силу 


фиксированности~$\chi$. Элементарные свойства оператора $\xi_\Lambda$ 


определяет 


 


\begin{lem} 


Пусть $\Lambda^{\#}_\sigma$ и $X_\sigma$ --- пространства 


$\Lambda^{\#}$ и $X$, снабженные слабыми топологиями 


$\sigma(\Lambda^{\#},\Lambda^{\# *})$, $\sigma(X,X^*)$ соответственно, 


где $\Lambda^{\# *}$ и $X^*$ --- пространства, топологически сопряженные 


к пространствам $\Lambda^{\#}$ и $X$. Тогда {\em a)} 


$\xi_\Lambda\in\cal L(\Lambda^{\#},X);$ {\em б)} $\xi_\Lambda\in\cal 


L(\Lambda^{\#},X_\sigma);$ {\em в)} $\xi_\Lambda\in\cal 


L(\Lambda^{\#}_\sigma,X_\sigma)$. 


\end{lem} 


 


{\bf Доказательство.} Включение $\xi_\Lambda\in \cal 


L(\Lambda^{\#},X)$ очевидно, б) --- его прямое следствие, поэтому 


установим лишь включение в). Для каждого $x^*\in X^*$ в силу б) 


отображение $\lambda\to\la x^*,\xi_\Lambda(\lambda)\ra:\Lambda^{\#}\to 


R$ линейно и непрерывно, где $\la\cdot,\cdot\ra$ означает каноническую 


билинейную форму, устанавливающую двойственность между $X$ и $X^*$. 


Таким образом, $\la x^*,\xi_\Lambda(\cdot)\ra\in\Lambda^{\# *}$ и, 


следовательно, для любого открытого в $R$ множества $Q$ справедливы 


включения $x^{*-1}[Q]\subset \sigma(X,X^*)$ и $\la 


x^*,\xi_\Lambda(\cdot)\ra^{-1}[Q]\subset 


\sigma(\Lambda^{\#},\Lambda^{\# *})$. Отсюда делаем заключение о 


справедливости утверждения в). 


 


Опишем геометрию образов оператора $\xi_\Lambda$ со ``стандартных'' 


множеств из $\Lambda^{\#}$. Следуя общепринятой терминологии, 


множество $\{\lambda\in \Lambda_k:\|\lambda_0-\lambda\|_{\nu,k}\ge 


r\}$, $k\in[p',\infty]$, назовем замкнутым $r$-шаром в точке 


$\lambda_0\in\Lambda_k$ и обозначим $\ov S^k_r(\lambda_0)$. В следующей 


теореме конкретизируется топологическая структура множества откликов 


$\xi$-модели на воздействия сигнальных функций из $\ov 


S^k_r(\lambda_0)\subset \Lambda^{\#}$. 


 


\begin{thm} 


Область значений оператора $\xi_\Lambda:\Lambda^{\#}_{p'}\to X$, 


$p'\in[1,\infty)$, сепарабельна, множество $\xi_\Lambda[\ov 


S^k_r(\lambda_0)]$, $k\in[p',\infty]$, $p'\in(1,\infty)$, компактно и 


метризуемо в $X_\sigma$. 


\end{thm} 


 


\begin{cor} 


Пусть $p'=\infty$, тогда $\xi_\Lambda[\ov S^\infty_r(\lambda_0)]$ --- 


компакт в $X_{\sigma}$, при этом он будет метризуем, если $X$ 


сепарабельно. 


\end{cor} 


 


{\bf 4. Структура максимального идентификационного базиса.} Для 


приложений общей теории $I$-процессов динамических объектов (1) большое 


значение имеют свойства пространственной формы максимального 


идентификационного базиса. 


 


Пусть $\chi\in H$ --- некоторый характеристический элемент $I$-процесса 


и $\Lambda^{\#}$ --- максимальное семейство его сигнальных функций. 


Обозначим через $\Omega^{\#}$ идентификационный базис вида 


$\{\lambda\chi\in H;\ \lambda\in\Lambda^{\#}\}$. Ясно, что 


$\Omega^{\#}$ --- наибольший $I$-базис в семействе $I$-базисов, 


отвечающих $\chi$. Условимся через $\frak J^{\#}_\Omega$ обозначать 


идентификацию на $\cal L(H,X)$, индуцированную $I$-базисом 


$\Omega^{\#}$. 


 


На двух теоремах, следующих ниже, основан стандартный подход к задачам 


описания нетривиальных максимальных подмножеств, идентифицируемых в 


$(\cal L(H,X),\frak J^{\#}_\Omega)$: так, например, чтобы подтвердить 


факт существования таких подмножеств, достаточно (если $\Omega^{\#}$ --- 


линейное подмножество) показать, что замыкание $\Omega^{\#}$ в $H$ 


является собственным подпространством $H$, откуда в силу теоремы 2 


следует, что $I$-базис $\Omega^{\#}$ неполный и, значит, пространство 


$(\cal L(H,X),\frak J^{\#}_\Omega)$ содержит собственные идентифицируемые 


подпространства, максимальные в упорядочении относительно 


теоретико-множественного включения (более подробно этот вопрос 


рассмотрим в разделе 5). 


 


\begin{thm} 


$\Omega^{\#}_{p'}$ --- собственное сепарабельное подпространство в 


$H_{p'}$, а при $p'>1$ идентификационный базис $\Omega^{\#}_{p'}$ --- 


сепарабельное равномерно выпуклое $($и, следовательно, рефлексивное$)$ 


банахово пространство в себе $($с нормой, индуцированной 


$\|\cdot\|_{H_{p'}})$. 


\end{thm} 


 


\begin{note} 


a) Здесь и далее подпространством называем замкнутое линейное 


множество в $H_{p'}$. б) Известно \cite{16}, что равномерно выпуклые 


пространства образуют весьма частный подкласс класса рефлексивных 


пространств. 


\end{note} 


 


{\bf Доказательство.} Легко проверить, что $0\in\Omega^{\#}$ и 


$\alpha\Omega^{\#}+\beta\Omega^{\#}\subset\Omega^{\#}$ 


($\alpha,\beta\in R$), откуда следует линейность подмножества 


$\Omega^{\#}\subset H$. Чтобы доказать $\ov\Omega^{\#}\ne H$ построим 


элемент $h\in H$, $h\notin\ov\Omega^{\#}$. Пусть 


$\lambda_0\in\Lambda^{\#}$ и 


$\Lambda_\lambda\hateq\{\alpha\lambda_0:\alpha\in R\}$. Ниже 


понадобится понятие расстояния от точки $\lambda^*$ до множества 


$\Lambda_\lambda$. Как и в эвклидовом пространстве, понимаем под этим 


величину $\rho(\lambda^*,\Lambda_\lambda)= 


\inf\limits_{\lambda\in\Lambda_\lambda}\|\lambda^*- 


\lambda\|_{\Lambda^{\#}}$. Ясно, что 


$\rho(\lambda^*,\Lambda_\lambda)=0$ равносильно 


$\lambda^*\in\Lambda_\lambda$. Далее, в соответствии с леммой 2 


(\cite{13}, с.\,129) для любого $\varepsilon>0$ существует сигнальная 


функция $\lambda_1\in\Lambda^{\#}$ такая, что 


$\rho(\lambda_1,\Lambda_\lambda)>1-\varepsilon$. Тогда в качестве 


элемента $h$ можно взять $(\lambda_0x,\lambda_1u)\in H$. Чтобы 


обосновать замкнутость $\Omega^{\#}_{p'}$ в $H_{p'}$ достаточно для 


тройки $(H_{p'},\Omega^{\#}_{p'},\Lambda^{\#}_{p'})$ подтвердить свойство 


полноты $\Omega^{\#}_{p'}$ в топологии, индуцированной исходной топологией 


в $H_{p'}$. С этой целью рассмотрим оператор 


$\varkappa_{p'}:\Lambda^{\#}_{p'}\to\Omega^{\#}_{p'}$ вида 


$\varkappa_{p'}(\lambda)(t)\hateq\lambda(t)\chi(t)$, $t\in T$, который 


является линейной биекцией. В силу теоремы 4 для всех $\lambda\in 


\Lambda^{\#}_{p'}$\; 


$\|\varkappa_{p'}(\lambda)\|_{H_{p'}}=\|\lambda\|_{\nu,p'}$. Таким 


образом, оператор $\varkappa_{p'}$ --- линейная изометрия, отображающая 


банахово пространство $\Lambda^{\#}_{p'}$ на $\Omega^{\#}_{p'}$. Отсюда 


с очевидностью следует полнота, сепарабельность и при $p'>1$ 


равномерная выпуклость \cite{17} линейного множества $\Omega^{\#}_{p'}$ 


в топологии, индуцированной нормой $\|\cdot\|_{H_{p'}}$ (при этом 


учтено, что $\nu$ --- мера со счетным базисом и, следовательно, 


пространство $\Lambda^{\#}_{p'}$ в силу п.\,a) теоремы 4 сепарабельно). 


 


\begin{cor} 


$\Lambda^{\#}_{p'}$ и $\Omega^{\#}_{p'}$ линейно изометричны. 


\end{cor} 


 


(Данное следствие полезно при анализе свойства ОЛД-расширения (см., 


напр., доказательство теоремы 2 из \cite{9}).) 


 


Для тройки $(H_\infty,\Omega^{\#}_\infty,\Lambda^{\#}_\infty)$ простая 


переформулировка теоремы 7 неверна. Чтобы устранить препятствия, 


возникающие на этом пути, рассмотрим одно вспомогательное построение. 


Назовем $\nu$-существенной нижней гранью функции 


$\psi\hateq(\omega_1,\omega_2)\in L_\infty(T,\mu,X)\times 


L_\infty(T,\mu,Y)$ наибольшую из констант $c$, для которых 


$\max(\|\omega_1(t)\|_X,\|\omega_2(t)\|_Y)\ge c$ имеет место 


$\nu$-почти всюду на $T$. Обозначим ее через 


$\nu\text{-}\essinf\|\psi\|_{X,Y}$. 


 


\begin{thm} 


Линейное множество $\Omega^{\#}_\infty$ замкнуто в $H_\infty$ тогда и 


только тогда, когда\linebreak $\nu\text{-}\essinf\|\chi\|_{X,Y}\ne0$. 


\end{thm} 


 


Быть или не быть элементу $\nu\text{-}\essinf\|\chi\|_{X,Y}$ обратимым 


--- чисто алгебраическое обстоятельство. С другой стороны, замкнутость 


$\Omega^{\#}_\infty$ зависит от метрических свойств $H_\infty$. В этом 


состоит одно из неожиданных заключений теоремы 8 --- она устанавливает 


совпадение величин совершенно различного происхождения. 


 


Линейная изометрия $\varkappa_{p'}$, построенная при доказательстве 


теоремы 7, а также равенство $\xi_\Lambda=\xi\circ\varkappa_{p'}$, 


делают очевидным 


 


\begin{assrt} 


При $p'\in[1,\infty)$ теорема $6$ и лемма $2$ будут справедливыми, если 


в них $\Lambda^{\#}$ заменить на $\Omega^{\#}$, а $\xi_\Lambda$ --- на 


$\xi$. 


\end{assrt} 


 


Бесконечномерный вариант (аналог) теоремы 5 из \cite{2} обобщает 


 


\begin{thm} 


Пусть $\Lambda_\sigma$ --- семейство всех характеристических функций 


$\sigma$-алгебры $\frak S_\nu$ и пусть для любого $S\in\frak S_\nu$ 


такого, что $\nu(S)>0$, система функций $\{\la x^*,\Dot x(\cdot)\mid 


S\ra:x^*\in X^{\#}\}$, где $X^{\#}$ --- некоторое тотальное на $X$ 


подмножество из $X^*$, не является фундаментальной в $L_1(S,\nu,R)$. 


Тогда $\xi_\Lambda[\ov 


S^\infty_1(0)]=\xi_\Lambda[\Lambda_\sigma-\Lambda_\sigma]$. 


\end{thm} 


 


\begin{note} 


a) $\la x^*,\Dot x(\cdot)\mid S\ra\in L_1(S,\nu,R)$ следует из п.\,a) 


следствия 2; б) так как мера $\nu$ неатомическая, то посылка теоремы 9 


заведомо имеет место, если пространство $X$ конечномерно. 


\end{note} 


 


{\bf Доказательство.} Из $\Lambda_\sigma-\Lambda_\sigma\subset \ov 


S^\infty_1(0)$ следует 


$\xi_\Lambda[\Lambda_\sigma-\Lambda_\sigma]\subset \xi_\Lambda[\ov 


S^\infty_1(0)]$. Чтобы доказать обратное включение, зафиксируем 


какую-нибудь точку $z\in\xi_\Lambda[\ov S^\infty_1(0)]$, $z\ne0$ (при 


$z=0$ дальнейшие выкладки тривиальны) и положим 


$\Lambda_z\hateq\{\lambda\in\ov 


S^\infty_1(0):\xi_\Lambda(\lambda)=z\}$. Достаточно показать, что 


$\Lambda_z$ содержит некоторую функцию $\lambda=\chi_{S'}-\chi_{S''}$, 


$S',S''\in\frak S_\nu$, $S'\cap S''=\emptyset$, $\nu(S'\cup S'')\ne0$. 


Так как оператор $\xi_\Lambda(\cdot)$ слабо непрерывен на 


$\Lambda^{\#}_\infty$, то множество $\Lambda_z$ выпукло и 


$(\ast)$-слабо компактно. По теореме Крейна--Мильмана $\Lambda_z$ имеет 


крайнюю точку. Покажем, что функция, отвечающая этой точке, имеет вид 


$\chi_{S'}-\chi_{S''}$. 


 


Пусть $\lambda_0\in\Lambda_z$ и не является функцией вида 


$\chi_{S'}-\chi_{S''}$. Тогда найдутся такое $S\in\frak S_\nu$ и число 


$\varepsilon\in(0,1)$, что $\nu(S)>0$ и 


$-1+\varepsilon\le\lambda_0\le1-\varepsilon$ на $S$. Положим 


$Z\hateq\chi_s\Lambda^{\#}_\infty$. Если ядро оператора 


$\xi_\Lambda\mid Z$ нетривиально, то найдется такой элемент 


$\lambda^*\ne0$ в подпространстве $Z$, что $\xi_\Lambda\mid 


Z(\lambda^*)=0$ и $-\varepsilon<\lambda^*<\varepsilon$. Ясно, что 


$(\lambda_0+\lambda^*),\,(\lambda_0-\lambda^*)\in\Lambda_z$ и, 


следовательно, $\lambda_0$ не является крайней точкой множества 


$\Lambda_z$. Поэтому для завершения доказательства теоремы покажем, что 


$\ker\xi_\Lambda\mid Z(0)\ne\{0\}$. 


 


В соответствии с выражением (3) оператор 


$\xi_\Lambda:\Lambda^{\#}_\infty\to X$ в аналитической форме выглядит 


так: $\xi_\Lambda(\lambda)=\int\limits_T\lambda(t)\Dot x(t)\mu(dt)$. 


Следовательно, $\lambda\in\ker\xi_\Lambda\mid Z$ в том и только том 


случае, если $\int\limits_S\lambda(t)\la x^*,\Dot x(t)\ra\mu dt=0$ для любого 


$x^*\in X^{\#}$. В свою очередь, последнее утверждение справедливо при 


и только при условии, что $\forall x^*\in X^{\#}$ имеет место 


$\int\limits_S\lambda(t)\la x^*,\Dot x(t)\ra\nu(dt)=0$. С другой 


стороны, поскольку множество $\{\la x^*,\Dot x(\cdot)\mid S\ra:x^*\in 


X^{\#}\}$ не является фундаментальным в $L_1(S,\nu,R)$, то согласно 


теореме 3 найдется такая функция $\lambda\in L_\infty(S,\nu,R)$, что 


$\lambda\ne0$ и $\forall x^*\in X^{\#}$ будет 


$\int\limits_S\lambda(t)\la x^*,\Dot x(t)\ra\nu(dt)=0$. 


 


\begin{cor} 


В силу следствия 3 множество $\xi_\Lambda[\Lambda_\sigma-\Lambda_\sigma]$ 


выпукло, компактно и если $X^{\#}$ счетно, то метризуемо в $X_\sigma$. 


\end{cor} 


 


В заключение раздела остановимся кратко на проблеме существования в 


$\Omega^{\#}$ минимального $I$-базиса, порождающего ту же идентификацию 


$\frak J^{\#}_\Omega$ на $\cal L(H,X)$. Ясно, что если такой 


минимальный $I$-базис существует, то в соответствии с теоремой 2 при 


сохранении свойства минимальности в $\Omega^{\#}$ он должен 


образовывать множество, фундаментальное в $\Omega^{\#}$ (в 


предположении, что топология в $\Omega^{\#}$ индуцирована исходной 


топологией в $H$). Очевидно, что в качестве такого множества мог бы 


выступать базис банахова пространства $\Omega^{\#}$. С другой стороны, 


проблема существования такого базиса (см. краткую справку в \cite{13}, 


с.\,514) имеет в силу следствия 4 положительное решение при $H=H_{p'}$, 


$p'\in[1,\infty)$. 


\medskip 


 


{\bf 5. Геометрия максимальных идентифицируемых подмножеств в $(\cal 


L(H,X),\frak J^{\#}_\Omega)$.} На заданном отрезке времени практик, как 


правило, имеет лишь одну траекторию (см., напр., \cite{6}), 


характеризующую модель системы, что для некоторых классов динамических 


объектов (в частности, с уравнением состояния (1)) должно принципиально 


ограничивать такое их свойство как идентифицируемость. В этих случаях 


понятие максимального идентифицируемого подмножества возникает 


естественным образом, а исследование его тополого-алгебраической 


структуры вызывает неформальный интерес. 


 


Будем рассматривать задачу: для идентификационного пространства $(\cal 


L(H,X),\frak J^{\#}_\Omega)$ построить семейство $\cal M$ всех 


максимальных идентифицируемых подмножеств в $\cal L(H,X)$, при которых 


структурно описать $\cal M$ как семейство


топологическиx векторных пространств (ТВП) с топологиями поточечной 


сходимости на $\Omega^{\#}$ (постановка B из \cite{1}). Впервые 


подобная задача была поставлена авторами в статье \cite{2}, содержащей 


результаты для конечномерного объекта (1). В действительности основные 


черты теории максимальных идентификационных подмножеств и сила 


функционально-аналитического подхода, лежащего в ее основании, в полной 


мере проявляются уже в конечномерном случае. Однако сама природа 


методов и результатов делает распространение их на ТВП настолько 


естественным, что можно идти ради этого на некоторые технические 


усложнения и детали. 


 


Пусть $\pi_\omega:\cal L(H,X)\to X$ --- отображение вычисления в точке 


$\omega\in\Omega^{\#}$ и $V^{\#}\subset\cal L(H,X)\times\cal L(H,X)$ 


--- идентификационный признак, равный 


$\bigcap\limits_{\omega\in\Omega^{\#}}\pi^{-1}_\omega\circ\pi_\omega$ 


(то, что $\bigcap\limits_{\omega\in\Omega^{\#}}\pi^{-1}_\omega\circ 


\pi_\omega\ne\Lambda_L$, гарантировано теоремами 2, 7). Обозначим через 


$\frak D^{\#}$ разбиение $\cal L(H,X)$, отвечающее отношению 


эквивалентности $V^{\#}$. В соответствии с теоремой 7 из \cite{1} имеет 


место следующее утверждение: $M\in\cal M\Longleftrightarrow 


M=\bigcup\limits_{D\in\frak D^{\#}}\{\xi_D\}\, \&\, \xi_D\in D\in\frak 


D^{\#}$. Следовательно, для каждого $M\in\cal M$ существует линейный 


взаимно однозначный оператор $\zeta_M:M\to\cal L(H,X)/ 


\bigcap\limits_{\omega\in\Omega^{\#}}\ker\pi_\omega$ вида 


$\zeta_M(\xi_D)\hateq V^{\#}(\xi_D)=D$, $D\in\frak D^{\#}$. Далее, в 


силу теоремы 8 из \cite{1} все элементы в $\cal M$ идентификационно 


изоморфны. Поэтому, если $(M,\frak J^M)$ и $(N,\frak J^N)$ --- 


идентификационные подпространства пространств $(\cal L(H,X),\frak 


J^{\#}_\Omega)$ и $M,N\in\cal M$, то оператор 


$\zeta^{-1}_M\circ\zeta_M:M\to N$ --- идентификационный изоморфизм. 


Определим на элементах из $\cal M$ алгебраические структуры, а именно, 


введем бинарную операцию $\oplus$ сложения и операцию $\odot$ умножения 


на скаляры по правилам 


\begin{alignat*}{2} 


\xi_1\oplus\xi_2&\hateq \zeta^{-1}_M(\zeta_M(\xi_1)+ 


\zeta_M(\xi_2)), &\quad &\xi_1,\xi_2\in M;\\ 


%


\alpha\odot\xi_3&\hateq \zeta^{-1}_M(\alpha\zeta_M(\xi_3)), 


&\quad & \alpha\in R,\ \; \xi_3\in M, 


\end{alignat*} 


при этом $\zeta^{-1}_M[\bigcap\limits_{\omega\in\Omega^{\#}} 


\ker\pi_\omega]\in M$ назовем нулевым элементом пространства 


$(M,\oplus,\odot)$ и обозначим через $0_M$. Ясно, что операция $\oplus$ 


ассоциативна и коммутативна, а в сочетании с $\odot$ выполняются 


дистрибутивность и ассоциативность умножения. 


 


Ходом предыдущих рассуждений по существу было доказано 


 


\begin{assrt} 


Тройка $(M,\oplus,\odot)$ --- векторное пространство, алгебраически 


изоморфное фактор-пространству $\cal 


L(H,X)/\bigcap\limits_{\omega\in\Omega^{\#}}\ker\pi_\omega$. 


\end{assrt} 


 


В большинстве случаев, когда рассматривается конкретное векторное 


пространство, в нем уже имеется некоторая ``естественная'' сходимость, 


которая определяет топологию этого пространства. То, что 


$(M,\oplus,\odot)$ не исключение, показывает 





\begin{thm} 


Пусть $\cal T_M$ --- топологическая структура на $M\in\cal M$, 


представляющая топологию поточечной сходимости на $\Omega^{\#}$. Тогда 


$(M,\oplus,\odot,\cal T_M)$ --- ТВП. 


\end{thm} 


 


(Так как большинство интересных теорем анализа содержат условия 


``хаусдорфовости'' в качестве одного из предположений, то ниже включаем 


вторую аксиому отделимости в число аксиом ТВП.) 


\medskip 


 


{\bf Доказательство.} Обозначим через $\cal L_S(H,X)$ пространство 


$\cal L(H,X)$ с топологией поточечной сходимости на $H$, которое 


является ТВП. Заметим, что отображение $\pi_\omega:\cal L_S(H,X)\to X$ 


--- линейный непрерывный оператор и, следовательно, пространство 


$\bigcap\limits_{\omega\in\Omega^{\#}}\ker\pi_\omega$ замкнуто в $\cal 


L_S(H,X)$. Таким образом, на $\cal 


L_S(H,X)/\bigcap\limits_{\omega\in\Omega^{\#}}\ker\pi_\omega$ можно 


корректно определить фактор-топологию, превращающую 


$\cal L_S(H,X)/\bigcap\limits_{\omega\in\Omega^{\#}}\ker\pi_\omega$ в 


ТВП (см., напр., теорему 1.41 из \cite{15}). Далее, множества $U^{\omega 


Q}_M\hateq\{v\in M:v(\omega)\in Q\}$ и $U^{\omega Q}_{\cal L}\hateq 


\{D\in\frak D^{\#}: \exists\xi_D\in D\,\& \, \xi_D(\omega)\in Q\}$, где 


$\omega\in\Omega^{\#}$ и $Q$ --- открытая область в $X$, образуют 


предбазы соответственно в $\cal T_M$ и фактор-топологии на 


$\cal L_S(H,X)/\bigcap\limits_{\omega\in\Omega^{\#}}\ker\pi_\omega$. 


Нетрудно показать, что $\zeta^{-1}_M(D)(\omega)=\xi_D(\omega)$, где 


$\xi_D\in M\, \&\, \xi_D\in D$ и, следовательно, 


$\zeta^{-1}_M[U^{\omega Q}_{\cal L}]=U^{\omega Q}_M$, для любых $\omega$ 


и $Q$, определенных выше. Таким образом, $\zeta_M$ --- линейный 


гомеоморфизм $(M,\oplus,\odot,\cal T_M)$ на 


$\cal L_S(H,X)/\bigcap\limits_{\omega\in\Omega^{\#}}\ker\pi_\omega$. 


 


\begin{note} 


a) Теорема останется справедливой при замене исходной топологии в 


банаховом пространстве $X$ на $\sigma\la X,X^*\ra$. б) Оператор 


$\zeta^{-1}_N\circ\zeta_M:M\to N$ ($M,N\in\cal M$) осуществляет 


линейный гомеоморфизм $(M,\oplus,\odot,\cal T_M)$ на 


$(N,\oplus,\odot,\cal T_N)$. 


\end{note} 


 


{\bf Заключение.} В работе развит геометрический подход к анализу 


свойств абстрактных идентификационных процессов в классе линейных 


динамических систем управления с уравнениями состояния в общем 


банаховом пространстве. Подход основан на идеях и математических 


конструкциях предложенной в \cite{1} теоретико-множественной 


аксиоматизации общей теории $I$-процессов. 


 


Изложенная в данной статье методология является достаточно гибким и 


эффективным средством анализа формальных идентификационных процессов в 


классе объектов (1). Она позволяет во многих практических случаях, не 


касаясь алгоритмической природы используемого метода идентификации, 


отвечать на ``фундаментальные'' вопросы; так, например, свойство 


замкнутости максимального $I$-базиса одним из основных следствий 


устанавливает, что принципиально невозможно по реализации одной 


траектории движения однозначно определить математическую модель 


динамического объекта (1) (даже в конечномерном случае). Таким 


образом, в отличие от других известных методов анализа свойств 


идентифицируемости систем вида (1) и, как правило, предлагающих узкие 


интерпретации идентифицируемости, изложенный в работе подход позволяет 


не только решать некоторые частные задачи системного анализа, но и дает 


универсальную возможность с общих позиций подходить к исследованию 


свойств $I$-процессов, при этом формулировать утверждения о таких 


свойствах с помощью единой системы понятий и терминов. 


 


В качестве приоритетных направлений дальнейшего продвижения общей 


теории $I$-процессов в классе динамических объектов с уравнениями 


состояния в общем банаховом пространстве по мнению авторов могут стать: 


 


--- развитие теории в предположении, когда наблюдается выход 


$y(t)=C(t)x(t)$, где $C$ --- оператор наблюдателя; 


 


--- изучение идентификационных инвариантов на максимальных 


идентифицируемых фактор-пространствах $(A,B)$- и $\xi$-моделей; 


 


-- построение теории сильных $(A,B)$-моделей (см. \cite{3}--\cite{5}, 


\cite{7}--\cite{9}, \cite{18}), связанное с принципом 


$I$-эквивалентности \cite{3}; 


 


-- обоснование $I$-процессов распределенных систем \cite{19}. 


 


 


\begin{thebibliography}{99} 


 


\bibitem{1} 


Данеев А.В., Русанов В.А. {\em К аксиоматической теории идентификации 


динамических систем}. I. {\em Основные структуры} // Автоматика и 


телемеханика. -- 1994. -- \No\,8. -- С.\,126--136. 


\bibitem{2} 


Данеев А.В., Русанов В.А. {\em К аксиоматической теории идентификации 


динамических систем}. II. {\em Идентификация линейных систем} // 


Автоматика и телемеханика. -- 1994. -- \No\,9. -- С.\,120--133. 


\bibitem{3} 


Данеев А.В., Русанов В.А. {\em Об одной теореме существования сильной 


модели} // Автоматика и телемеханика. -- 1995. -- \No\,8. -- 


С.\,64--73. 


\bibitem{4} 


Vassilyev S.N., Rusanov V.A., Daneev A.V. {\em About automatic 


construction of mathematical models by methods of structural and 


parametrical identification} // CESA'96 Proc. Symp. on applied


mathematics and optimization. IMACS--IEE/SMC Multiconference. 


Computation engineering in systems applications. -- Lille. France, 


1996. -- P.\,27--31. 


\bibitem{5} 


Данеев А.В., Русанов В.А. {\em Геометрические характеристики свойств 


существования конечномерных $(A,B)$-моделей в задачах 


структурно-параметрической идентификации} // Автоматика и телемеханика. 


-- 1999. -- \No\,1. -- С.\,3--8. 


\bibitem{6} 


Данеев А.В., Куменко А.Е., Русанов В.А. {\em Задача спектральной 


идентификации математической модели линейной динамической системы 


управления ЛА} // Изв. вузов. Авиационная техника. -- 1999. -- \No\,1. -- 


С.\,20--24. 


\bibitem{7} 


Данеев А.В., Русанов В.А. {\em Порядковые характеристики свойств 


существования сильных линейных конечномерных дифференциальных моделей} 


// Дифференц. уравнения. -- 1999. -- Т.\,35. -- \No\,1. -- С.\,43--50. 


\bibitem{8} 


Lakeyev A.V., Rusanov V.A. {\em On existence of linear nonstationary 


partially realizing models} // CESA'98 Proc. Symp. on applied 


mathematics and optimization. IMACS--IEE/SMC Multiconference. 


Computational engineering in systems applications. -- 1998. -- 


P.\,201--205. 


\bibitem{9} 


Данеев А.В., Русанов В.А. {\em Об одном классе сильных дифференциальных 


моделей над счетным множеством динамических процессов конечного 


характера} // Изв. вузов. Математика. -- 2000. -- \No\,2. -- 


С.\,32--40. 


\bibitem{10} 


Калман Р., Фалб П., Арбиб М. {\em Очерки по математической теории 


систем}. -- М.: Мир, 1971. -- 400 с. 


\bibitem{11} 


Ahmed N.U. {\em Optimization and identification of systems governed by 


evolution equation on Banach space}. -- New York: John Wiley and Sons, 


1988. -- 188 p. 


\bibitem{12} 


Иосида К. {\em Функциональный анализ}. -- М.: Мир, 1967. -- 624 с. 


\bibitem{13} 


Канторович Л.В., Акилов Г.П. {\em Функциональный анализ}. -- М.: Наука, 


1977. -- 742 с. 


\bibitem{14} 


Варга Дж. {\em Оптимальные управления дифференциальными и 


функциональными уравнениями}. -- М.: Наука, 1977. -- 624 с. 


\bibitem{15} 


Рудин У. {\em Функциональный анализ}. -- М.: Мир, 1975. -- 444 с. 


\bibitem{16} 


Pettis B.J. {\em A proof that every uniformly convex space is 


reflexive} // Duke Math. J. -- 1939. -- V.\,5. -- P.\,249--253. 


\bibitem{17} 


Clarkson J.A. {\em Uniformly convex spaces} // Trans. Amer. Math. Soc. 


-- 1936. -- V.\,40. -- P.\,396--414. 


\bibitem{18} 


Русанов В.А., Данеев А.В., Дмитриев А.В., Мартьянов В.И. {\em 


Аналитический подход к теории структурной идентификации$:$ 


характеризация линейных дифференциальных моделей управления на основе 


оператора Рэлея--Ритца} // Тр. Международн. конф. ``Идентификация 


систем и задачи управления'' (SISPRO'2000). -- М., 26--28 сентября, 


2000. -- С.\,525--536. 


\bibitem{19} 


Rusanov V.A., Daneev A.V., Dmitriev A.V. {\em The spectral analysis of 


$I$-processes in the class of mixed problems for linear models of 


normal-hyperbolic type} // Proc. 14-th World Congress of IFAC, July 


5--9, 1999. -- Beijing, China. -- V.\,H. -- P.\,409--414. 


 


 


\end{thebibliography} 


 


\vskip 2cm 


 


\post{\it Институт динамики}{\it Поступила} 


\post{\it систем и теории управления}{22.11.1999} 


\post{\it Сибирского отделения}{} 


\post{\it Российской Академии наук}{} 


\smallskip 


 


\post{\it Иркутский государственный}{} 


\post{\it технический университет}{} 


 


\label{end} 


\end{document} 


